\documentclass[11pt]{ctexart}
\usepackage[a4paper,margin=1in]{geometry}
\usepackage{graphicx}
\usepackage{xcolor}
\usepackage{hyperref}
\definecolor{refgray}{gray}{0.45}
\title{\textbf{市场最新观点汇总}\\[0.4em]{\large 市场主线：全球资产定价正从需求驱动转向供给侧结构性重塑，AI硬件与工业投资超级周期成为新引擎，而中国市场的流动性收缩与消费断层构成主要下行风险。}}
\author{KC桌面 自动生成}
\date{260528}
\begin{document}
\maketitle
\section*{一页摘要}
\begin{itemize}
  \item 全球石油边际成本下降并非需求疲软，而是生产端效率改善，长期油价锚点可能被低估。
  \item 汽车半导体涨价周期已至，AI服务器功耗溢出与上游代工成本传导是结构性推手。
  \item 日本股市重估的核心是硬件与实物资产的结构性优势，而非AI泡沫。
  \item 中国央行正以创纪录速度抽走中期流动性，收益率曲线陡峭化是必然结果。
  \item 美债外国官方卖盘是汇率干预的副产品，而非结构性减持。
  \item 美国消费韧性依赖一次性减税红利，收入断层与信心脱钩构成下半年风险。
  \item 中国住房可负担性已重置至十年前水平，月供能力改善被系统性低估。
  \item 亚洲非科技出口复苏的宽度被低估，工业投资超级周期正在形成。
  \item 中国供给侧行政收缩（铝、煤炭、钢铁）正在重塑成本曲线与定价逻辑。
  \item AI电力需求正从短期主题转化为长期产能规划，发电设备供给紧平衡将持续。
\end{itemize}
\section{宏观与利率：流动性范式切换与消费断层}
\textbf{全球宏观环境正经历流动性管理范式的切换，中国央行“收长放短”重塑收益率曲线，美国消费韧性面临收入断层与信心脱钩的结构性挑战。}

\begin{itemize}
  \item 中国央行通过买断式逆回购创纪录地净回笼中期流动性，累计规模超2.5万亿元，未来三个月到期量高达5.7万亿元，收益率曲线陡峭化是必然结果。
  \item 美国消费者正经历“结构性分裂”：资产负债表强劲，但收入增长和消费信心已跌至历史低位，实际消费支出增速下半年将低于共识预期。
  \item 美债外国官方卖盘的核心驱动力是美元走强引发的汇率干预，而非资产偏好转移，一旦冲突缓和卖盘大概率逆转。
  \item 中国自由流动性指标跌至2025年5月以来新低，M1增长乏力与PPI加速上行共同侵蚀货币供给的真实购买力。
  \item 美国大型银行GSIB附加资本在现行规则下中位数环比上升50个基点，监管资本成本的重新定价将影响资本分配与股东回报。
\end{itemize}
\begin{figure}[htbp]
\centering
\includegraphics[width=0.88\linewidth]{figures/fig_117.jpg}
\caption{Figure 1 - Figure 1: Largest medium-term liquidity tightening on record RMBbn}
\end{figure}
\begin{figure}[htbp]
\centering
\includegraphics[width=0.88\linewidth]{figures/fig_118.jpg}
\caption{Figure 2 - Figure 2: RMB5.7tn upcoming maturities for MLF and ORR in next 3 months Maturities}
\end{figure}
\begin{figure}[htbp]
\centering
\includegraphics[width=0.88\linewidth]{figures/fig_129.jpg}
\caption{Exhibit 1 - Exhibit 2: We Expect Softer Spending Growth in 2026H2}
\end{figure}
\begin{figure}[htbp]
\centering
\includegraphics[width=0.88\linewidth]{figures/fig_130.jpg}
\caption{Exhibit 2 - Exhibit 2: We Expect Softer Spending Growth in 2026H2}
\end{figure}
\begin{figure}[htbp]
\centering
\includegraphics[width=0.88\linewidth]{figures/fig_136.jpg}
\caption{Exhibit 7 - Exhibit 7: Consumer Sentiment Has Pulled Back Sharply Since the Start of the War in Iran}
\end{figure}
\begin{figure}[htbp]
\centering
\includegraphics[width=0.88\linewidth]{figures/fig_121.jpg}
\caption{Exhibit 1 - Exhibit 1: Foreign official holdings of Treasuries have declined as a share of the market Foreign Ownership of US Treasury Securities}
\end{figure}
\begin{figure}[htbp]
\centering
\includegraphics[width=0.88\linewidth]{figures/fig_126.jpg}
\caption{Exhibit 6 - Exhibit 6: USD performance has historically been a key explanatory factor for foreign official demand, reflecting policy motivations}
\end{figure}
\begin{figure}[htbp]
\centering
\includegraphics[width=0.88\linewidth]{figures/fig_160.jpg}
\caption{Exhibit 1 - Exhibit 1: MS China Free Liquidity Indicator vs. MSCI China YoY \% – Free liquidity dropped further in April, renewing the lowest level since May 2025; we expect it to stay negative through 2026}
\end{figure}
\begin{figure}[htbp]
\centering
\includegraphics[width=0.88\linewidth]{figures/fig_289.jpg}
\caption{Exhibit 1 - Exhibit 1: GSIB scores increased q/q across the group in 1Q26 under current rules, with the biggest increase at JPM Fed GSIB Score, Current Rules MS \& CO. LLC \# Manan Gosalia}
\end{figure}
\begin{figure}[htbp]
\centering
\includegraphics[width=0.88\linewidth]{figures/fig_290.jpg}
\caption{Exhibit 4 - Exhibit 4: GSIB scores increased q/q across the board in 1Q26 under current rules, with the biggest increase at JPM Fed GSIB Score, Current Rules}
\end{figure}
{\small\color{refgray}\textbf{References:} [R004] 市场低估的不是降息与否，而是央行正在以创纪录速度抽走中期流动性; [R005] 市场真正低估的不是资金紧张，而是央行流动性管理范式的切换; [R007] 美国消费的真正风险不在支出放缓，而在收入断层与信心脱钩; [R006] 美债市场真正担心的不是抛售，而是抛售背后的动机被误读; [R010] 市场真正低估的不是需求，而是流动性的结构性收缩; [R032] 市场真正低估的不是需求，而是GSIB附加资本的重新定价}

\section{AI/算力与半导体：从算力竞赛到网络架构重构}
\textbf{AI基础设施的投资逻辑正从算力芯片转向网络架构与电力设备，半导体供给侧的永久性收缩与架构创新正在重塑定价权。}

\begin{itemize}
  \item 微软数据中心容量四年内翻两番，资本开支远超当前变现速度，但建设超前于收入意味着未来收入预测可能被系统性低估。
  \item 谷歌TPU 8系列首次将训练与推理网络架构彻底分离，带动光学交换、高速光模块及ARM架构CPU需求的系统性增长。
  \item 华为Tau Scaling Law标志着从“几何缩放”到“时间缩放”的转变，先进封装设备商与EDA厂商的议价能力正在系统性上升。
  \item DDR4供需缺口从14\%扩大至19-20\%，美光设备迁移与SK海力士产能腾挪导致供给永久性收缩，涨价周期正在路上。
  \item PC半导体市场正经历供给侧重构，功耗优势正在改写竞争格局，份额迁移成为2027年增长的核心变量。
  \item 中国半导体设备自给率突破21\%，沉积和刻蚀环节已形成双轮驱动，存储厂IPO加速产能扩张创造增量需求。
\end{itemize}
\begin{figure}[htbp]
\centering
\includegraphics[width=0.88\linewidth]{figures/fig_274.jpg}
\caption{Exhibit 1 - Exhibit 1: We Estimate Microsoft's Revenue per MW Declines as the Company Accelerates DC Capacity Expansion, Supporting Upside to Our Revenue Estimates Over Time Microsoft Cloud Revenue per MW of Datacenter Capacity}
\end{figure}
\begin{figure}[htbp]
\centering
\includegraphics[width=0.88\linewidth]{figures/fig_275.jpg}
\caption{Exhibit 2 - Exhibit 2: Capex-Implied Azure AI Revenue Potential is Well Ahead of Our Azure AI Forecast}
\end{figure}
\begin{figure}[htbp]
\centering
\includegraphics[width=0.88\linewidth]{figures/fig_252.jpg}
\caption{Figure 3 - Figure 3. Share Price Movement – by Sectors © 2026 Citi Inc. No redistribution without Citi's written permission. Figure 4. Share Price Movement – 1 Day}
\end{figure}
\begin{figure}[htbp]
\centering
\includegraphics[width=0.88\linewidth]{figures/fig_253.jpg}
\caption{Figure 5 - 1-Day Avg Px Move © 2026 Citi Inc. No redistribution without Citi's written permission. Figure 5. Share Price Movement – 1 Week}
\end{figure}
\begin{figure}[htbp]
\centering
\includegraphics[width=0.88\linewidth]{figures/fig_259.jpg}
\caption{Exhibit 1 - Exhibit 1: Order of preference: GC memory Exhibit 2: NOR flash demand and supply growth rates}
\end{figure}
\begin{figure}[htbp]
\centering
\includegraphics[width=0.88\linewidth]{figures/fig_260.jpg}
\caption{Exhibit 6 - Exhibit 6: Quarterly supply breakdown (mn Gb)}
\end{figure}
\begin{figure}[htbp]
\centering
\includegraphics[width=0.88\linewidth]{figures/fig_261.jpg}
\caption{Exhibit 8 - Exhibit 8: DDR4 8Gb (1Gx8) pricing chart}
\end{figure}
\begin{figure}[htbp]
\centering
\includegraphics[width=0.88\linewidth]{figures/fig_263.jpg}
\caption{Exhibit 5 - Exhibit 5: Quarterly oversupply/undersupply ratio vs. Nanya and Winbond pricing Q/Q change}
\end{figure}
\begin{figure}[htbp]
\centering
\includegraphics[width=0.88\linewidth]{figures/fig_288.jpg}
\caption{Exhibit 4 - Exhibit 4: Parade: Historical P/E}
\end{figure}
\begin{figure}[htbp]
\centering
\includegraphics[width=0.88\linewidth]{figures/fig_225.jpg}
\caption{Exhibit 1 - EXHIBIT 1: Global WFE: China demand expected to be weaker while non-China could still go up in 2025 Global \& China Wafer Fab Equipment (WFE) TAM USD bn}
\end{figure}
\begin{figure}[htbp]
\centering
\includegraphics[width=0.88\linewidth]{figures/fig_226.jpg}
\caption{EXHIBIT 2 - EXHIBIT 2: China WFE: Acceleration in domestic substitution continues, expect the self-sufficiency to reach 26\% by 2026 China Wafer Fab Equipment (WFE) TAM and Domestic Share USD bn}
\end{figure}
\begin{figure}[htbp]
\centering
\includegraphics[width=0.88\linewidth]{figures/fig_227.jpg}
\caption{EXHIBIT 3 - EXHIBIT 3: In 2025, Chinese players kept growing in Deposition \& Dry Etch where U.S. players were traditionally stronger China WFE TAM, China supply \& Self-sufficiency by segment (2025)}
\end{figure}
\begin{figure}[htbp]
\centering
\includegraphics[width=0.88\linewidth]{figures/fig_276.jpg}
\caption{Exhibit 1 - Exhibit 1: China's WFE market – Our forecasts}
\end{figure}
\begin{figure}[htbp]
\centering
\includegraphics[width=0.88\linewidth]{figures/fig_277.jpg}
\caption{Exhibit 2 - Exhibit 2: Chinese WFE vendors is expected to account for 25\% of local WFE market share in 2026e}
\end{figure}
{\small\color{refgray}\textbf{References:} [R028] 微软的AI基建投入，市场低估的不是需求，而是变现的时间差; [R034] 市场真正低估的不是算力，而是网络架构的重新定价; [R022] 市场真正低估的不是华为的芯片，而是“时间缩放”对半导体定价权的重构; [R027] 市场低估的不是需求，而是供给侧的永久性收缩; [R031] PC半导体正在经历一轮被低估的供给侧重构; [R018] 中国半导体设备自给率突破20\%的真正含义：不是总量，是结构性拐点; [R029] 市场真正低估的不是需求，而是中国半导体设备供给侧的再定价}

\section{能源与大宗：供给侧收缩与成本曲线重塑}
\textbf{大宗商品市场的定价逻辑正从需求驱动转向供给侧意外收缩与成本曲线重塑，中国行政干预与安全监管成为关键变量。}

\begin{itemize}
  \item 全球石油边际成本降至69美元/桶，但市场误读为天花板，实际是效率改善，2026年可能反弹至77美元/桶。
  \item 山西煤矿事故导致约1.25亿吨原煤产能暂停，安全生产三年行动方案正从阶段性运动转向制度性约束。
  \item 环保督查组进驻电解铝主产区，确认超产现象，可能影响40-50万吨产能，库存堆积与价格飙升并存。
  \item 中国尿素出口窗口重启，设定价格底线高于多数国际基准，是一次有管理的供给侧再定价。
  \item 钢铁市场成本曲线因焦煤供给中断与海运运费飙升而重塑，出口量已攀升至历史区间上沿。
  \item 全球电力设备进入超级周期，AI数据中心、电网现代化与新能源并网三股力量叠加，备用电源订单已排至2028年。
\end{itemize}
\begin{figure}[htbp]
\centering
\includegraphics[width=0.88\linewidth]{figures/fig_001.jpg}
\caption{EXHIBIT 2 - EXHIBIT 2: The Marginal Cost of Oil for Top 50 decreased by 2\% to US\textbackslash{}\$69/bbl in 2025}
\end{figure}
\begin{figure}[htbp]
\centering
\includegraphics[width=0.88\linewidth]{figures/fig_002.jpg}
\caption{Exhibit 4 - EXHIBIT 4: 2025 Global non-OPEC ex FSU Marginal Cost Curve}
\end{figure}
\begin{figure}[htbp]
\centering
\includegraphics[width=0.88\linewidth]{figures/fig_028.jpg}
\caption{EXHIBIT 32 - EXHIBIT 32: We view that US\textbackslash{}\$75/bbl is a reasonable estimate for the long term price of oil in equity valuation, which is above the current 60M forward strip of US\textbackslash{}\$70/bbl}
\end{figure}
\begin{figure}[htbp]
\centering
\includegraphics[width=0.88\linewidth]{figures/fig_151.jpg}
\caption{Figure 1 - Figure 1: Total China steel output 1,026Mt annualised run rate: daily CISA steel output for 10-days ended 20 May: -1\% vs previous (+1\% YoY), in line with seasonal trend}
\end{figure}
\begin{figure}[htbp]
\centering
\includegraphics[width=0.88\linewidth]{figures/fig_152.jpg}
\caption{Figure 2 - Figure 2: FOB iron ore price from Australia to China at \textbackslash{}\textasciitilde{}\textbackslash{}\$89/t, flat vs end of Feb}
\end{figure}
\begin{figure}[htbp]
\centering
\includegraphics[width=0.88\linewidth]{figures/fig_154.jpg}
\caption{Figure 4 - Figure 4: Iron ore bulk shipping costs from Australia, Brazil and South Africa to China are all \$+45\textbackslash{}\%\$ since start of the Iran conflict LHS: Iron ore price \textbackslash{}\$/mt; RHS: Shipping cost \textbackslash{}\$/mt}
\end{figure}
\begin{figure}[htbp]
\centering
\includegraphics[width=0.88\linewidth]{figures/fig_254.jpg}
\caption{Exhibit 1 - Exhibit 1: Recent year coal accidents impacted domestic spot prices as safety measures constrained supply in the subsequent months Exhibit 2: The start of a 3-year safety act (安全生产治本攻坚三年行动方案) following the accidents in 2H23 led t}
\end{figure}
\begin{figure}[htbp]
\centering
\includegraphics[width=0.88\linewidth]{figures/fig_255.jpg}
\caption{Exhibit 1 - Exhibit 1: Recent year coal accidents impacted domestic spot prices as safety measures constrained supply in the subsequent months Exhibit 2: The start of a 3-year safety act (安全生产治本攻坚三年行动方案) following the accidents in 2H23 led t}
\end{figure}
\begin{figure}[htbp]
\centering
\includegraphics[width=0.88\linewidth]{figures/fig_256.jpg}
\caption{Exhibit 3 - Exhibit 3: Monthly china coal imports declined in 2026 YTD by an average of 43mnt versus 2025A Exhibit 4: Coal inventories at the top 6 IPPs were below normal seasonalities as of May-2026}
\end{figure}
\begin{figure}[htbp]
\centering
\includegraphics[width=0.88\linewidth]{figures/fig_257.jpg}
\caption{Exhibit 3 - Exhibit 3: Monthly china coal imports declined in 2026 YTD by an average of 43mnt versus 2025A Exhibit 4: Coal inventories at the top 6 IPPs were below normal seasonalities as of May-2026}
\end{figure}
\begin{figure}[htbp]
\centering
\includegraphics[width=0.88\linewidth]{figures/fig_190.jpg}
\caption{Exhibit 1 - Exhibit 1: Historical Chinese Urea Exports Chinese Urea Exports}
\end{figure}
\begin{figure}[htbp]
\centering
\includegraphics[width=0.88\linewidth]{figures/fig_191.jpg}
\caption{Exhibit 2 - Exhibit 2: Chinese Urea Export Seasonality}
\end{figure}
{\small\color{refgray}\textbf{References:} [R001] 市场真正低估的不是油价，而是供给侧的再定价空间; [R009] 钢铁市场的真正叙事不在需求端，而在供给侧的意外收缩与成本曲线的重塑; [R023] 煤炭市场的真正风险不是事故本身，而是安全监管的系统性收紧; [R012] 铝行业真正超预期的不是需求，而是供给侧的行政收缩正在加速; [R017] 铝价暴涨，市场真正定价的不是库存，而是中国供给侧的“自我阉割”; [R013] 中国尿素出口窗口重启：市场真正低估的不是需求，而是供给侧的再定价; [R026] AI 电力需求不是短期主题，它正在重塑全球发电设备供给格局}

\section{地产与消费：可负担性改善与结构性分化}
\textbf{中国住房可负担性已重置至十年前水平，月供能力改善被低估；美国消费韧性依赖一次性政策红利，收入断层构成风险。}

\begin{itemize}
  \item 中国一线城市房价收入比从21.5降至16.3，月供金额较2021年高点下降约42\%，利率下行与房价修正叠加产生乘数效应。
  \item 上海4月二手房周交易量创2016年以来新高，低总价房源占比稳定在64\%，二手房流动性恢复正在激活“卖旧买新”链条。
  \item 美国4月零售数据强劲主要受OBBBA减税带来的约1400亿美元一次性收入提振，下半年通胀拖累将显现。
  \item 美国消费者信心跌至历史低位，收入增长与消费信心脱钩，实际可支配收入同比增速仅0.4\%。
  \item 中国量贩零食赛道补贴战并非2024年翻版，而是有纪律的“定点防御”，盈利影响有限。
  \item 爱马仕二手溢价倍数从1.2倍回升至1.5倍，稀缺性定价权未被破坏，市场低估其供给侧的定价韧性。
\end{itemize}
\begin{figure}[htbp]
\centering
\includegraphics[width=0.88\linewidth]{figures/fig_137.jpg}
\caption{Figure 1 - Figure 1: Broad-based improvement in home price affordability on par with the level a decade ago}
\end{figure}
\begin{figure}[htbp]
\centering
\includegraphics[width=0.88\linewidth]{figures/fig_138.jpg}
\caption{Figure 2 - Figure 2: Improved mortgage payment ability}
\end{figure}
\begin{figure}[htbp]
\centering
\includegraphics[width=0.88\linewidth]{figures/fig_140.jpg}
\caption{Figure 6 - Figure 6: Shanghai April secondary transactions hit a new high since 2016}
\end{figure}
\begin{figure}[htbp]
\centering
\includegraphics[width=0.88\linewidth]{figures/fig_141.jpg}
\caption{Figure 7 - Figure 7: Home prices correct 11-25\% from the peak as of March 2026}
\end{figure}
\begin{figure}[htbp]
\centering
\includegraphics[width=0.88\linewidth]{figures/fig_142.jpg}
\caption{Figure 8 - Figure 8: Mortgage rate down 2.6ppt from the last peak}
\end{figure}
\begin{figure}[htbp]
\centering
\includegraphics[width=0.88\linewidth]{figures/fig_129.jpg}
\caption{Exhibit 1 - Exhibit 2: We Expect Softer Spending Growth in 2026H2}
\end{figure}
\begin{figure}[htbp]
\centering
\includegraphics[width=0.88\linewidth]{figures/fig_130.jpg}
\caption{Exhibit 2 - Exhibit 2: We Expect Softer Spending Growth in 2026H2}
\end{figure}
\begin{figure}[htbp]
\centering
\includegraphics[width=0.88\linewidth]{figures/fig_132.jpg}
\caption{Exhibit 4 - Exhibit 4: We Forecast Only 1.3\% Real income Growth in 2026 Due to Energy Price Headwinds, With Significant Underperformance Among Lower-Income Households and Weak Cashflow Growth (Adjusted for Timing of OBBBA Tax Cuts) in 2026H2}
\end{figure}
\begin{figure}[htbp]
\centering
\includegraphics[width=0.88\linewidth]{figures/fig_136.jpg}
\caption{Exhibit 7 - Exhibit 7: Consumer Sentiment Has Pulled Back Sharply Since the Start of the War in Iran}
\end{figure}
\begin{figure}[htbp]
\centering
\includegraphics[width=0.88\linewidth]{figures/fig_235.jpg}
\caption{EXHIBIT 1 - EXHIBIT 1: Our ‘core’ resale price tracker, which only includes data on the Birkin 25/30 and Kelly 25/28 Hermès Birkin \& Kelly Resale Premium (volume weighted average)}
\end{figure}
\begin{figure}[htbp]
\centering
\includegraphics[width=0.88\linewidth]{figures/fig_236.jpg}
\caption{EXHIBIT 2 - EXHIBIT 2: Our expanded tracker now includes data on the Mini Kelly II and Kelly Pochette Hermès Birkin, Kelly, Mini Kelly \& Kelly Pochette Resale Premium (volume weighted average)}
\end{figure}
{\small\color{refgray}\textbf{References:} [R008] 市场真正低估的不是需求，而是月供能力的历史性重置; [R007] 美国消费的真正风险不在支出放缓，而在收入断层与信心脱钩; [R024] 市场真正担心的不是补贴战，而是补贴战之后谁还能涨; [R019] 爱马仕的二手价格正在释放一个比财报更重要的信号}

\section{区域市场：亚洲工业投资超级周期与日本硬件优势}
\textbf{亚洲正在进入2000年代中期以来最强劲的工业投资周期，非科技出口复苏的宽度被低估；日本股市重估的核心是硬件与实物资产的结构性优势。}

\begin{itemize}
  \item 亚洲非科技出口在过去半年以年化27\%的速度增长，覆盖资本品、中间品和消费品多个领域，AI、能源转型与国防支出是三大引擎。
  \item 亚洲资本开支预计从11万亿美元增长至2030年的16万亿美元，年复合增速7\%，是2023-25年增速的三倍。
  \item 日本企业利润结构发生质变，电气设备与银行业贡献主要增量，保守指引后往往出现阶梯式上涨。
  \item 日本机床对中国出口订单同比增长57\%，由一般机械和电机精密机械驱动，全球制造业资本开支进入分化加速阶段。
  \item 中国新能源车市场品牌间订单表现极端分化，新车型脉冲效应取代稳态增长成为核心竞争力。
\end{itemize}
\begin{figure}[htbp]
\centering
\includegraphics[width=0.88\linewidth]{figures/fig_161.jpg}
\caption{Exhibit 1 - Exhibit 1: Asia's gross fixed investment to rise to US\textbackslash{}\$16trn by 2030 Asia nominal investment (US\$trn)}
\end{figure}
\begin{figure}[htbp]
\centering
\includegraphics[width=0.88\linewidth]{figures/fig_162.jpg}
\caption{Exhibit 2 - Exhibit 2: The rise in capex is to be driven by new structural demand drivers – AI and AI infrastructure spending, energy and energy transition, and defense spending Asia nominal investment in high growth sectors (US\$trn)}
\end{figure}
\begin{figure}[htbp]
\centering
\includegraphics[width=0.88\linewidth]{figures/fig_165.jpg}
\caption{Exhibit 5 - Exhibit 5: A sharp rise in Asia's non-tech exports since October 2025}
\end{figure}
\begin{figure}[htbp]
\centering
\includegraphics[width=0.88\linewidth]{figures/fig_166.jpg}
\caption{Exhibit 6 - Exhibit 6: Non-tech exports for early reporters in Asia grew at \$16\textbackslash{}\%\$ Y in April Asia exports ex semiconductors ex fuels (\%Y, 3mma)}
\end{figure}
\begin{figure}[htbp]
\centering
\includegraphics[width=0.88\linewidth]{figures/fig_167.jpg}
\caption{Exhibit 7 - Exhibit 7: The non-semis export recovery has been broad-based across end-use product segments}
\end{figure}
\begin{figure}[htbp]
\centering
\includegraphics[width=0.88\linewidth]{figures/fig_168.jpg}
\caption{Exhibit 8 - Exhibit 8: Capital goods and intermediate goods drove the first leg of the recovery, while consumer goods exports have also joined in the recovery}
\end{figure}
\begin{figure}[htbp]
\centering
\includegraphics[width=0.88\linewidth]{figures/fig_192.jpg}
\caption{Exhibit 1 - Exhibit 1: Machine tool orders: Exports by region}
\end{figure}
\begin{figure}[htbp]
\centering
\includegraphics[width=0.88\linewidth]{figures/fig_193.jpg}
\caption{Exhibit 1 - Exhibit 1: A sharp rise in Asia's non tech exports since October 2025}
\end{figure}
\begin{figure}[htbp]
\centering
\includegraphics[width=0.88\linewidth]{figures/fig_194.jpg}
\caption{Exhibit 2 - Exhibit 2: A sharp rise in Asia's non tech exports since October 2025}
\end{figure}
\begin{figure}[htbp]
\centering
\includegraphics[width=0.88\linewidth]{figures/fig_195.jpg}
\caption{Exhibit 3 - Exhibit 3: Non-tech exports for early reporters in Asia grew at \$15\textbackslash{}\%\$ Y in April}
\end{figure}
\begin{figure}[htbp]
\centering
\includegraphics[width=0.88\linewidth]{figures/fig_200.jpg}
\caption{Exhibit 8 - Exhibit 8: On an annualized basis, non-tech exports for Asia's early reporting economies have grown at \$27\textbackslash{}\%\$ since October through April}
\end{figure}
\begin{figure}[htbp]
\centering
\includegraphics[width=0.88\linewidth]{figures/fig_201.jpg}
\caption{Exhibit 9 - Exhibit 9: On a YoY three-month trailing basis, Asia's non-tech exports growth has accelerated to \$14\textbackslash{}\%\$ Y Exports ex semiconductors ex fuels (\%Y 3mma, as of April-26 or latest)}
\end{figure}
{\small\color{refgray}\textbf{References:} [R011] 亚洲市场真正被低估的不是AI，而是工业投资超级周期的宽度; [R015] 亚洲的真正主线不是半导体，而是被低估的非科技出口复苏; [R003] 日本股市的真正重估：不是AI泡沫，而是硬件与实物资产的结构性优势; [R014] 中国订单重回历史峰值，但市场真正该关注的不是总量，而是结构; [R021] 中国新能源车市真正的问题不是需求疲软，而是供给侧的“分化陷阱”; [R030] 中国电动车市场真正的博弈，已经从“卖多少”转向“靠什么卖”}

\section{风险偏好与市场情绪：仓位出清与催化剂等待}
\textbf{中国资产仓位出清已近尾声，市场焦点从“是否配置”转向“什么因素能触发趋势逆转”，AI硬件热潮的可持续性与政策风险是核心变量。}

\begin{itemize}
  \item MSCI中国在新兴市场指数中的权重从31\%下降至21\%，仓位调整已非常充分，进一步下跌动能衰竭。
  \item 投资者讨论焦点已转向趋势逆转的催化剂：利率上升、大型IPO抽走流动性、AI资本开支能否持续盈利。
  \item 中国创新药在ASCO的“量变”已确认，94项口头报告创历史新高，但少数头部公司在特定靶点上的“局部定价权”尚未被充分定价。
  \item NSCLC一线治疗正从PD-1单药+化疗进入双抗、ADC和新型免疫疗法驱动的差异化竞争阶段。
  \item 市场对中国资产的悲观预期已充分定价，但趋势逆转需要明确的催化剂，当前处于“共识形成之后、催化剂出现之前”的灰色地带。
\end{itemize}
\begin{figure}[htbp]
\centering
\includegraphics[width=0.88\linewidth]{figures/fig_246.jpg}
\caption{Exhibit 1 - Exhibit 1: China underperformed YTD Performance of US and major Asian markets YTD}
\end{figure}
\begin{figure}[htbp]
\centering
\includegraphics[width=0.88\linewidth]{figures/fig_247.jpg}
\caption{Exhibit 3 - Exhibit 3: China's weighting fell to below Taiwan and Korea Major EM markets' weighting in MSCI EM BofA GLOBAL RESEARCH Exhibit 5: Semis and energy sectors outperformed YTD MSCI China sector performance YTD}
\end{figure}
\begin{figure}[htbp]
\centering
\includegraphics[width=0.88\linewidth]{figures/fig_248.jpg}
\caption{Exhibit 3 - Exhibit 3: China's weighting fell to below Taiwan and Korea Major EM markets' weighting in MSCI EM BofA GLOBAL RESEARCH Exhibit 5: Semis and energy sectors outperformed YTD MSCI China sector performance YTD}
\end{figure}
\begin{figure}[htbp]
\centering
\includegraphics[width=0.88\linewidth]{figures/fig_250.jpg}
\caption{Exhibit 4 - Exhibit 4: 2026E EPS growth was cut from 11\% in Jan to 5\% in May Consensus forecast for MSCI China's earning growth}
\end{figure}
{\small\color{refgray}\textbf{References:} [R020] 市场低估的不是中国资产的吸引力，而是“趋势逆转”的催化剂正在酝酿; [R033] 中国创新药在ASCO的“量变”已确认，但“质变”的定价权尚未被市场充分消化; [R025] NSCLC一线治疗格局正在被重新定义，但真正的赢家尚未完全浮现}

\section*{结语}
综合来看，当前市场主线清晰：供给侧的结构性重塑（石油边际成本、半导体产能、中国行政限产）与AI驱动的工业投资超级周期是两大核心叙事，而中国流动性收缩与美国消费断层构成主要下行风险。投资者需从总量思维转向结构思维，关注那些被共识低估的供给侧变量与区域分化机会。
\vfill
{\small\color{refgray}Personal reading notes and learning share only. Not investment advice.}
\end{document}
